\chapter{Conclusion and Future Scope}

\section{Conclusion}

UltraFabric-Vision set out to build a real-time, unsupervised fabric-defect detection system by fusing three complementary anomaly detectors---PatchCore density estimation, DINO self-supervised attention, and Vision Transformer Autoencoder reconstruction---into a single calibrated decision, to minimize inference latency, and to industrialize the result as a deployable service. These three objectives structured the entire project.

The objectives were met through a combination of diagnostic and engineering work. A preprocessing defect that had silently destroyed accuracy---inference-only enhancement operators absent at fitting time---was identified and corrected by enforcing byte-level train--inference parity. The three detectors' incomparable raw scores were standardized to a common z-score scale using statistics estimated on defect-free data, yielding a single interpretable decision variable and a portable threshold. Both nearest-neighbour searches were migrated to a single GPU kernel and augmented with mixed precision and single-pass attention, and the calibrated engine was wrapped in a REST/WebSocket service, an embeddable SDK, a batch CLI, and a GPU container.

Quantitatively, the calibrated ensemble achieves perfect ranking on the evaluation benchmark (AUROC $=1.000$, $F_1 = 1.000$) with a wide decision margin: defect-free cloth clusters near a z-score of $0.17$ while defects reach a mean of $80.2$, against a threshold of $17.31$. The three pathways are quantitatively diverse (defect-score correlations of $0.71$--$0.72$ between DINO and the other detectors), and the full service is verified end-to-end by an automated test suite. Real-time throughput was confirmed on an NVIDIA RTX~4050 laptop GPU: the full ensemble runs at $37.3$~ms per frame ($26.8$~FPS, fp16) in only $641$~MB of memory, achieving high accuracy and low latency at the same time. The system was further demonstrated on the batch-video workflow, correctly localizing defects along a batch by position and generating per-batch quality-control records. The corrected preprocessing and score calibration together transform an unreliable prototype into a system whose decisions are accurate, reproducible, and real-time.

\section{Future Scope}

The principal limitation of the present study is that its quantitative results are obtained on a controlled synthetic dataset; validation on real production fabric, across weaves and materials, is the immediate next step. Further directions include: on-hardware GPU benchmarking to confirm the projected line-speed latency; online adaptation of the memory banks to accommodate gradual process drift without full retraining; temporal fusion across consecutive frames to suppress transient false alarms; distillation of the three-model ensemble into a single lightweight student for edge deployment; and extension from binary detection to multi-class defect categorization, for which the autoencoder's bimodal response already provides a starting signal.

\section{Learning Outcomes of the Project}
\begin{itemize}
\item Understanding of unsupervised anomaly detection and how memory-bank, self-supervised, and reconstruction-based methods differ in their inductive biases.
\item Practical insight into Vision Transformers and self-attention as feature extractors for industrial inspection.
\item The importance of train--inference consistency: how a subtle preprocessing mismatch can silently destroy the accuracy of distance-based models.
\item Techniques for calibrating and fusing heterogeneous model outputs into a single, interpretable decision variable.
\item GPU acceleration of machine-learning inference, including batched distance computation, mixed precision, and warm-up.
\item Engineering a research prototype into a deployable, configurable, and tested service, including REST/WebSocket APIs, an embeddable SDK, and containerized deployment.
\end{itemize}
